%% Julia Steinberg — Curriculum Vitae
%% Compiled with XeLaTeX or pdfLaTeX
%% Fonts: requires lmodern; swap  for any preferred font

\documentclass[11pt, letterpaper]{article}

% ── Packages ──────────────────────────────────────────────────────────────────
\usepackage[margin=1in, top=0.85in, bottom=0.85in]{geometry}
\usepackage{draftwatermark}
\SetWatermarkScale{3}
\SetWatermarkText{Draft}
\usepackage[T1]{fontenc}
\usepackage[utf8]{inputenc}
\usepackage{amsmath}
\usepackage{enumitem}
\usepackage{hyperref}
\usepackage{xcolor}
\usepackage{titlesec}
\usepackage{tabularx}
\usepackage{longtable}
\usepackage{array}
\usepackage{parskip}

% ── Colors ────────────────────────────────────────────────────────────────────
\definecolor{accentblue}{RGB}{26, 86, 160}   % section header color
\definecolor{darkgray}{RGB}{55, 65, 81}      % body text
\definecolor{midgray}{RGB}{107, 114, 128}    % secondary info

% ── Hyperlink style ───────────────────────────────────────────────────────────
\hypersetup{
  colorlinks = true,
  urlcolor   = accentblue,
  linkcolor  = accentblue,
  pdfauthor  = {Julia Steinberg},
  pdftitle   = {Julia Steinberg — CV}
}

% ── Section formatting ────────────────────────────────────────────────────────
\titleformat{\section}
  {\large\bfseries\color{accentblue}}   % format
  {}                                     % label
  {0em}                                  % sep
  {}                                     % before-code
  [\vspace{-4pt}\textcolor{accentblue}{\rule{\linewidth}{0.6pt}}\vspace{2pt}]

\titlespacing{\section}{0pt}{14pt}{4pt}

% ── List style ────────────────────────────────────────────────────────────────
\setlist[itemize]{
  leftmargin = 1.4em,
  itemsep    = 1pt,
  topsep     = 2pt,
  parsep     = 0pt,
  label      = \textbullet
}

% ── Helper: dated entry ───────────────────────────────────────────────────────
% Usage: \entry{Title · Organization, Location}{Date(s)}
\newcommand{\entry}[2]{%
  \noindent
  \begin{tabularx}{\linewidth}{@{}X@{\hspace{4pt}}r@{}}
    \textbf{#1} & \textcolor{midgray}{\small #2}
  \end{tabularx}\vspace{1pt}%
}

% ── Helper: publication entry ─────────────────────────────────────────────────
\newcommand{\pub}[2]{%
  \noindent\hangindent=2em\hangafter=1
  \textbf{#1.}\enspace #2\par\vspace{3pt}%
}

% ── No page numbers ───────────────────────────────────────────────────────────
\pagestyle{empty}

% ══════════════════════════════════════════════════════════════════════════════
\begin{document}
% ══════════════════════════════════════════════════════════════════════════════

% ── Header ────────────────────────────────────────────────────────────────────
\begin{center}
  {\LARGE\bfseries Julia Steinberg}\\[5pt]
  {\small
    New York, NY\enspace·\enspace
    \href{mailto:jasteinberg24@gmail.com}{jasteinberg24@gmail.com}\enspace·\enspace
    +1\,(847)\,507--7468\\[2pt]
    \href{https://scholar.google.com}{Google Scholar}\enspace·\enspace
    \href{https://linkedin.com}{LinkedIn}
    \enspace·\enspace
    \href{https://github.com/jasteinberg}{GitHub}
    \enspace·\enspace
    \href{https://jasteinberg.github.io}{Website}
  }
\end{center}

\vspace{4pt}
\section{Background}

Harvard Physics PhD with a track record of research in quantum condensed matter and computational neuroscience, including work at the intersection of statistical physics, machine learning, and neural computation. Most recently, three years as a quantitative researcher in high-frequency trading, designing and running large-scale experiments to evaluate forecast models and trading strategies, and contributing to the deployment of these models into live trading systems.

% ── Technical Skills ──────────────────────────────────────────────────────────
\section{Technical Skills}

\noindent\textbf{Programming:} Python (NumPy, SciPy), C++, MATLAB\\
\textbf{Research areas:} Computational neuroscience, Statistical physics of learning and disordered systems, Quantum condensed matter

% ── Research Interests ────────────────────────────────────────────────────────
%\section{Research Interests}
%Computational neuroscience\enspace·\enspace
%Statistical physics of learning and disordered %systems\enspace·\enspace
%Associative memory\enspace·\enspace
%Representation of structured knowledge\enspace·\enspace
%Quantum condensed matter\enspace·\enspace

% ── Education ─────────────────────────────────────────────────────────────────
\section{Education}

\entry{Ph.D., Physics · Harvard University}{2020}
\begin{itemize}
  \item Thesis: \href{http://140.247.122.253/theses/Julia_Steinberg_thesis.pdf}{\textit{Universal Aspects of Quantum-Critical Dynamics In and Out of
        Equilibrium}} · Advisor: Subir Sachdev
\end{itemize}

\vspace{4pt}
\entry{A.M., Physics · Harvard University}{2019}

\vspace{4pt}
\entry{M.S., Physics · University of Pennsylvania}{2014}
\vspace{4pt}\entry{B.A., Physics · University of Pennsylvania}{2014}
\begin{itemize}
  \item Magna cum laude with distinction in physics · Minor in mathematics
\end{itemize}

% ── Experience ────────────────────────────────────────────────────────────────
\section{Experience}

\entry{Quantitative Researcher · Radix Trading LLC, New York, NY}{Mar 2023 – Mar 2026}
\begin{itemize}
  \item Researcher on the core quant team designing and running large-scale, multi-stage experiments for forecast model development and strategy optimization
  \item Responsible for deploying models into live production systems; contributed to
        shared C++ production codebase with commits to core
        trading binaries.
  \item Research on forecasting/latency tradeoffs and development of latency models for research simulation.
\end{itemize}

\vspace{6pt}
\entry{Associate Research Scholar / CPBF Fellow · Center for the Physics of Biological
Function, Princeton University, Princeton, NJ}{Aug 2020 – Mar 2023}
\begin{itemize}
  \item Independent postdoctoral researcher at the CUNY/Princeton Center for the
        Physics of Biological Function
  \item Developed a theoretical model of associative memory for structured knowledge
        (episodic memories, cognitive maps, semantic structures) — work bearing directly
        on how neural networks represent, store, and retrieve relational and hierarchical information.
  \item Periodic visiting researcher, Simons Center for Systems Biology, Institute for
        Advanced Study (2020–2022).
  \item Invited speaker: APS March Meeting, Les Houches School of Physics, van
        Vreeswijk Theoretical Neuroscience Seminar Series, Context and Episodic Memory
        Symposium, Flatiron Institute NeuroAI group, and others.
\end{itemize}

\vspace{6pt}
\entry{Postdoctoral Fellow, Swartz Program in Theoretical Neuroscience · Center for
Brain Science, Harvard University, Cambridge, MA}{Jun 2020 – Aug 2020}
\begin{itemize}
  \item Developed theories of learning and long-term memory using minimal neural network
        models, using tools and techniques from physics of disordered systems and statistical learning theory, combining analytical theory and numerical simulation.
\end{itemize}

\vspace{6pt}
\entry{Graduate Researcher · Department of Physics, Harvard University, Cambridge, MA}{Aug 2014 – May 2020}
\begin{itemize}
  \item Studied dynamical properties of strongly correlated quantum systems:
        relaxation and thermalization rates, transport coefficients, and onset of
        many-body quantum chaos, using tools from statistical physics, quantum field theory, and
        gauge-gravity duality .
  \item Swartz Program in Theoretical Neuroscience: extended graduate research into
        neural computation and learning theory.
\end{itemize}

\vspace{6pt}
\entry{Teaching Fellow · Harvard University, Cambridge, MA}{Sep 2019 – May 2020}
\begin{itemize}
  \item Physics 232 (Advanced Electromagnetism, grad), AM 226 (Neural Computation,
        grad), Physics 15A (Intro Mechanics, undergrad), AP 284 (Statistical Mechanics,
        grad).
  \item Derek Bok Certificate of Distinction in Teaching, Fall 2019.
\end{itemize}

\vspace{6pt}
\entry{Undergraduate Researcher · University of Pennsylvania \& University of Michigan}{2011 – 2014}
\begin{itemize}
  \item Rappe Group (Penn Chemistry): Electronic structure calculations using DFT to
        propose materials with novel band topology; collaboration with Kane and Mele
        (Penn Physics) on topological insulators as part of the Penn LRSM topological insulator seed.
  \item NSF Physics REU, University of Michigan (Summer 2013): Edge states in lattice
        models of 2D/3D topological semimetals.
  \item A.T.\ Johnson Group, Penn Physics (Summer 2011): Fabrication of graphene
        transistors using CVD and mechanical exfoliation.
\end{itemize}

% ── Publications ──────────────────────────────────────────────────────────────
\section{Peer-Reviewed Publications and Review Articles}

\textbf{\textcolor{accentblue}{Neural Networks, Learning Theory \& Statistical Physics of Machine Learning}}

\vspace{4pt}

\pub{10}{J.\ Steinberg, U.\ Adomaityt, A.\ Fachech, P.\ Mergny, D.\ Barbier, and
  R.\ Monasson. Replica method for computational problems with randomness: principles
  and illustrations. \textit{Journal of Statistical Mechanics: Theory and Experiment}
  2024\,(10), 104002.}

\pub{9}{J.\ Steinberg. Performance Capacity of a Complex Neural Network.
  \textit{Physics} 16, 108 (2023).}

\pub{8}{J.\ Steinberg and H.\ Sompolinsky. Associative memory of structured knowledge.
  \textit{Scientific Reports} 12, 21808 (2022). bioRxiv: 2022.02.22.481380.}

\pub{7}{J.\ Steinberg, M.\ Advani, and H.\ Sompolinsky. New role for circuit expansion
  for learning in neural networks. \textit{Physical Review E} 103, 022404 (2021).
  \textbf{[Editors' Suggestion]} arXiv:2008.08653.}

\vspace{6pt}
\textbf{\textcolor{accentblue}{Quantum Many-Body Physics \& Holography}}

\vspace{4pt}

\pub{6}{J.\ Steinberg and B.\ Swingle. Thermalization and many-body chaos in QED3.
  \textit{Physical Review D} 99, 076007 (2019). arXiv:1901.04984.}

\pub{5}{J.\ Steinberg, N.P.\ Armitage, F.H.L.\ Essler, and S.\ Sachdev. NMR
  relaxation in Ising spin chains. \textit{Physical Review B} 99, 035156 (2019).
  \textbf{[Editors' Suggestion]} arXiv:1810.11466.}

\pub{4}{A.\ Eberlein, V.\ Kasper, S.\ Sachdev, and J.\ Steinberg. Quantum quench of
  the Sachdev-Ye-Kitaev Model. \textit{Physical Review B} 96, 205123 (2017).
  arXiv:1706.07803.}

\pub{3}{A.\ Lucas and J.\ Steinberg. Charge diffusion and the butterfly effect in
  striped holographic matter. \textit{Journal of High Energy Physics} 10, 143 (2016).
  arXiv:1608.03286.}

\pub{2}{S.\ Chatterjee, Y.\ Qi, S.\ Sachdev, and J.\ Steinberg. Superconductivity
  from a confinement transition out of a FL* metal. \textit{Physical Review B} 94,
  024502 (2016). \textbf{[Editors' Suggestion]} arXiv:1603.03041.}

\pub{1}{J.\ Steinberg, S.M.\ Young, S.\ Zaheer, C.L.\ Kane, E.J.\ Mele, and
  A.M.\ Rappe. Bulk Dirac Points in Distorted Spinels. \textit{Physical Review
  Letters} 112, 036403 (2014). arXiv:1309.5967.}

% ── Invited Talks ─────────────────────────────────────────────────────────────
\section{Selected Invited Talks \& Seminars}
\vspace{-\parskip}
\begin{longtable}{@{}lp{5.2in}@{}}
  \textcolor{midgray}{Mar 2023} &
    Neural Networks: Insights from Physics to Models of Learning and Memory
    · \textit{Physics Colloquium, The College of New Jersey} \\[3pt]
  \textcolor{midgray}{Feb 2023} & Associative memory of structured knowledge
    · \href{https://statphysneuro.github.io}{\textit{Towards a Theory of Artificial and Biological Neural Networks Workshop, Les Houches School of Physics}} \\[3pt]
  \textcolor{midgray}{Oct 2022} &
    Associative memory of structured knowledge
    · \href{https://www.youtube.com/watch?v=rh_-Vh9CGbg}{\textit{van Vreeswijk Theoretical Neuroscience Seminar Series}} \\[3pt]
  \textcolor{midgray}{May 2022} &
    Associative memory of structured knowledge
    · \href{https://memory.psych.upenn.edu/CEMS_2022}{\textit{Context and Episodic Memory Symposium}} \\[3pt]
  \textcolor{midgray}{Mar 2022} &
    Associative memory of structured knowledge
    · \textit{APS March Meeting, Neural Systems III Focus Session} \\[3pt]
  \textcolor{midgray}{Mar 2022} &
    Associative memory of structured knowledge
    · \textit{NeuroAI Theory \& Manifolds Reading Group, Flatiron Institute} \\[3pt]
  \textcolor{midgray}{Sep 2021} &
    Associative memory of structured knowledge
    · \textit{Saxe Group, Gatsby Computational Neuroscience Unit /
    Sainsbury Wellcome Centre, UCL} \\[3pt]
  \textcolor{midgray}{Feb 2021} &
    Associative memory of structured knowledge
    · \textit{Kavli Institute for BioNano Science and Technology, Harvard} \\[3pt]
  \textcolor{midgray}{Jan 2021} &
    Associative memory of structured knowledge
    · \textit{Simons Center for Systems Biology, Institute for Advanced Study} \\[3pt]
  \textcolor{midgray}{Jan 2020} &
    A new role for sparse expansion in neural networks
    · \textit{Theory and Modeling of Living Systems, Emory University} \\[3pt]
  \textcolor{midgray}{Jan 2020} &
    A new role for sparse expansion in neural networks
    · \textit{Santa Fe Institute} \\[3pt]
  \textcolor{midgray}{Sep 2019} &
    A new role for sparse expansion in neural networks
    · \textit{Neuroscience Institute Seminar, NYU Langone Medical Center} \\
\end{longtable}

% ── Summer Schools ────────────────────────────────────────────────────────────
\section{Summer Schools \& Workshops}

\begin{tabularx}{\linewidth}{@{}lX@{}}
  \textcolor{midgray}{Jul 2022} &
    Statistical Physics of Machine Learning
    · \textit{Les Houches School of Physics} (Krzakala \& Zdeborova, organizers) \\[3pt]
  \textcolor{midgray}{Jul 2019} &
    Theoretical Biophysics
    · \textit{Boulder School 2019, University of Colorado Boulder} \\[3pt]
  \textcolor{midgray}{Dec 2016} &
    34th Jerusalem Winter School: New Horizons in Quantum Matter
    · \textit{Hebrew University of Jerusalem} \\
\end{tabularx}

% ── Honors ────────────────────────────────────────────────────────────────────
\section{Honors, Awards \& Fellowships}

\begin{tabularx}{\linewidth}{@{}lX@{}}
  \textcolor{midgray}{2020–2023} &
    CPBF Fellowship, CUNY/Princeton Center for the Physics of Biological Function \\[3pt]
  \textcolor{midgray}{2015–2020} &
    NSF Graduate Research Fellowship, National Science Foundation \\[3pt]
  \textcolor{midgray}{2014–2015} &
    E.M.\ Purcell Fellowship, Harvard University Physics \\[3pt]
  \textcolor{midgray}{2014} &
    James Mills Peirce Fellowship, Harvard Graduate School of Arts \& Sciences \\[3pt]
  \textcolor{midgray}{2014} &
    Rose Award for Outstanding Undergraduate Research, University of Pennsylvania \\[3pt]
    \textcolor{midgray}{Summer 2013} &
    NSF-REU University of Michigan Physics Summer Research Fellowship \\[3pt]
  \textcolor{midgray}{2013} &
    Roy and Diana Vagelos Science Challenge Award, University of Pennsylvania \\[3pt]
  \textcolor{midgray}{2012–2014} &
    University Scholar, University of Pennsylvania \\
    \textcolor{midgray}{Summer 2012} &
    NSF-REU Penn LRSM Summer Research Fellowship \\[3pt]
    \textcolor{midgray}{Summer 2011} &
    Penn Undergraduate Research Mentorship Program \\[3pt]
\end{tabularx}

% ── Teaching & Service ────────────────────────────────────────────────────────
\section{Service \& Outreach}

\begin{itemize}
  \item Referee: \textit{Physics Letters A}, \textit{Physical Review B},
        \textit{Physical Review Letters}, \textit{Physical Review X}
  \item Physics of Life Summer School Committee — Center for the Physics of
        Biological Function, Princeton/CUNY (2022)
  \item Mentor — Physics Undergraduate Mentoring Program, Princeton University
        (2021–2022)
  \item Mentor — UWIP Mentorship Families Program, Princeton Women in Physics (2020)
  \item GSMI Mentor — Científico Latino, supporting underrepresented graduate school
        applicants (2020–2021)
  \item Mentor — WiSTEM Mentorship Program, Harvard College Women's Center
        (2016–2020)
  \item Adopt-a-Physicist — American Physical Society, connecting with high school
        physics students (Fall 2017)
\end{itemize}

\end{document}
